\documentclass[10pt,a4paper]{amsart}
\usepackage[margin=19mm]{geometry}
\usepackage{amsmath,amssymb,mathtools}
\usepackage[scr=rsfs,cal=euler]{mathalfa}
\usepackage{xfrac}
\usepackage[unicode,hidelinks,bookmarks=false]{hyperref}
\newcommand{\ie}{i.e.\ }
\newcommand{\cf}{cf.\ }
\newcommand{\resp}{respectively\ }
\newcommand{\Subsection}{\subsection}
\providecommand{\nobreakdash}{\nobreak\hbox{-}}
\setlength{\emergencystretch}{2em}
\multlinegap0pt
\usepackage{bm}
\usepackage{bbm}
\usepackage{dsfont}
\usepackage{xspace}
\usepackage{cancel}
\usepackage{mathtools}
\usepackage{amsthm}
\makeatletter
\@ifundefined{lemma}{\newtheorem{lemma}{Lemma}}{}
\@ifundefined{proposition}{\newtheorem{proposition}{Proposition}}{}
\@ifundefined{theorem}{\newtheorem{theorem}{Theorem}}{}
\makeatother
\title[Weighted Product Inequalities for the Sine Function: A Gamma]{Weighted Product Inequalities for the Sine Function: A Gamma-Function Approach and Sharp Comparisons}
\author{Augustine L. Mahu, Beno\^it F. Sehba, Cecilia D. Williams}
\date{Source: arXiv:2604.13106v1 (CC BY 4.0)}
\begin{document}
\maketitle

\noindent\emph{Source and licence.} Weighted Product Inequalities for the Sine Function: A Gamma-Function Approach and Sharp Comparisons. Version arXiv:2604.13106v1 (CC BY 4.0), distributed under the Creative Commons Attribution 4.0 International licence, as recorded in the arXiv licence metadata for this exact version and shown on the abstract page https://arxiv.org/abs/2604.13106v1. Full rights notice: licenses/arxiv-2604.13106-CC-BY-4.0.txt.\par\smallskip

\noindent\textit{Editorial scope.} Counted unit: the Theorem whose source label is \texttt{thm:compare\_n} in the article (source0.tex lines 381--396, source label \texttt{thm:compare\_n}), delivered together with the article's own complete original proof of it (source0.tex lines 398--440, one \texttt{proof} environment of 43 lines). No mathematics is written, repaired or re-derived: statement and proof are the article's own text line for line. Required context retained uncounted: prop:main3 (lines 258--269); prop:main4 (lines 277--287); lem:sincompare (lines 334--344). Every definition, display and label of those lines is retained unchanged. Omitted from this package and not redistributed: 1--257, 270--276, 288--333, 345--380, 397--397, 441--1199. Nothing inside a retained excerpt is omitted or altered: every retained body is delivered exactly as the article's lines are written, so no omission receipt of any kind accompanies this package. Citations: the retained excerpts carry no citation command at all, so no bibliography entry of the article is transcribed and no reference is redistributed. Reference adaptation: none was needed and none was made; every reference inside the delivered unit resolves inside it, so no unresolved reference or citation is produced. Printed slips kept literal: no printed word, symbol, hypothesis or number of the counted statement or of its proof was altered. Layout and class: the article's own class is \texttt{amsart} with packages including inputenc, fontenc, amsmath, amssymb, amsthm, geometry, booktabs, color, xcolor, graphicx, hyperref, multicol, listings, caption; the wrapper is the lane's pinned \texttt{amsart} editorial wrapper at 10pt on A4, which restates the theorem-like environments the retained text needs and adds the article's own macro definitions that the retained text needs (none), with extra wrapper packages (bm, bbm, dsfont, xspace, cancel, mathtools). The delivered numbering is the unit's own. Assets: no retained span contains a figure environment, a graphics-inclusion command, a PDF-page inclusion, a table or a TikZ picture; no image file is copied into this package, the package declares no asset dependency and no image file is redistributed with it. Licence: version arXiv:2604.13106v1 of this article is distributed under the Creative Commons Attribution 4.0 International licence, as recorded in the arXiv licence metadata for this exact version and shown on the abstract page https://arxiv.org/abs/2604.13106v1. A full rights notice is delivered at licenses/arxiv-2604.13106-CC-BY-4.0.txt. Characters: every retained excerpt is pure ASCII, so the delivered engine, XeLaTeX with its default Latin Modern text font, reports no missing character.
\begin{proposition}\label{prop:main3}
Let $\lambda_1,\ldots,\lambda_n\in(0,1)$ and
$\lambda=\displaystyle\sum_{i=1}^n\lambda_i$.  Then
\begin{equation}\label{eq:propmain3}
  \Bigl(\prod_{i=1}^{n-1}\sin^{\lambda_{i+1}}(\pi\lambda_i)\Bigr)
  \sin^{\lambda_1}(\pi\lambda_n)
  \;\leq\;
  \sin^\lambda\!\left(\frac\pi\lambda
  \Bigl[\sum_{i=1}^{n-1}\lambda_{i+1}(1-\lambda_i)+\lambda_1\lambda_n
  \Bigr]\right).
\end{equation}
\end{proposition}

\begin{proposition}\label{prop:main4}
Under the hypotheses of Proposition~\ref{prop:main3},
\begin{equation}\label{eq:propmain4}
  \Bigl(\prod_{i=1}^{n-1}\sin^{\lambda_{i+1}}(\pi\lambda_i)\Bigr)
  \sin^{\lambda_1}(\pi\lambda_n)
  \;\leq\;
  \sin^\lambda\!\left(\frac\pi\lambda
  \Bigl[\sum_{i=1}^{n-1}\lambda_{i+1}\lambda_i+\lambda_1\lambda_n
  \Bigr]\right).
\end{equation}
\end{proposition}

\begin{lemma}\label{lem:sincompare}
Let $\lambda>0$ and $\sigma,\tau\in(0,\lambda)$.  Then
\[
  \sin\!\Bigl(\frac{\pi\sigma}{\lambda}\Bigr)
  \;\geq\;
  \sin\!\Bigl(\frac{\pi\tau}{\lambda}\Bigr)
  \quad\Longleftrightarrow\quad
  (\sigma-\tau)(\sigma+\tau-\lambda)\;\leq\;0,
\]
with equality if and only if $\sigma=\tau$ or $\sigma+\tau=\lambda$.
\end{lemma}

\begin{theorem}\label{thm:compare_n}
Let $\lambda_1,\ldots,\lambda_n\in(0,1)$, $\lambda=\sum_i\lambda_i$,
and let $S$, $T$ denote the upper bounds of
Propositions~\ref{prop:main3} and~\ref{prop:main4} respectively.
Define
\begin{equation}\label{eq:delta_n}
  \delta_n \;:=\; \sum_{i=1}^{n-1}\lambda_{i+1}(1-2\lambda_i).
\end{equation}
Then
\begin{equation}\label{eq:compare_n}
  S \;\geq\; T
  \quad\Longleftrightarrow\quad
  \delta_n\,(2\lambda_n-1) \;\leq\; 0,
\end{equation}
with equality $S=T$ if and only if $\delta_n=0$ or $\lambda_n=\tfrac12$.
\end{theorem}

\begin{proof}
Since $t\mapsto t^\lambda$ is increasing for $t>0$, $S\geq T$ if and only if
$\sin(\pi\sigma/\lambda)\geq\sin(\pi\tau/\lambda)$.
By Lemma~\ref{lem:sincompare}, this is equivalent to
$(\sigma-\tau)(\sigma+\tau-\lambda)\leq 0$.

\medskip
Let us compute $\sigma-\tau$ and $\sigma+\tau-\lambda$. We have
\begin{equation}\label{eq:sigmatau_n}
  \sigma-\tau
  =\sum_{i=1}^{n-1}\lambda_{i+1}(1-\lambda_i)-\sum_{i=1}^{n-1}\lambda_{i+1}\lambda_i
  =\sum_{i=1}^{n-1}\lambda_{i+1}(1-2\lambda_i)
  =\delta_n.
\end{equation}

\medskip
We also obtain
\begin{align*}
  \sigma+\tau
  &=\sum_{i=1}^{n-1}\lambda_{i+1}(1-\lambda_i)+\sum_{i=1}^{n-1}\lambda_{i+1}\lambda_i
   +2\lambda_1\lambda_n
  =\sum_{i=1}^{n-1}\lambda_{i+1}+2\lambda_1\lambda_n\\
  &=(\lambda-\lambda_1)+2\lambda_1\lambda_n,
\end{align*}
since $\sum_{i=2}^n\lambda_i=\lambda-\lambda_1$.  Therefore,
\begin{equation}\label{eq:sigmasum_n}
  \sigma+\tau-\lambda = \lambda_1(2\lambda_n-1).
\end{equation}

\medskip
Substituting \eqref{eq:sigmatau_n} and \eqref{eq:sigmasum_n} into
the criterion of Lemma~\ref{lem:sincompare} gives
\[
  \delta_n\cdot\lambda_1(2\lambda_n-1)\leq 0.
\]
Since $\lambda_1>0$, this is equivalent to $\delta_n(2\lambda_n-1)\leq 0$,
which is \eqref{eq:compare_n}.

The equality $S=T$ holds if and only if $\delta_n=0$ (i.e. $\sigma=\tau$) or
$2\lambda_n-1=0$, i.e. $\lambda_n=\tfrac12$.  

The proof is complete.
\end{proof}

\end{document}
